\subsection{文字列を返す}

これは\RobPikePractice{}からの古典的なバグです：

\begin{lstlisting}[style=customc]
#include <stdio.h>

char* amsg(int n, char* s)
{
        char buf[100];

        sprintf (buf, "error %d: %s\n", n, s) ;

        return buf;
};

int main()
{
        amsg ("%s\n", interim (1234, "something wrong!"));
};
\end{lstlisting}

これはクラッシュするでしょう。
まず、なぜかを理解しましょう。

これはamsg()が戻る前のスタック状態です：

\begin{lstlisting}
(lower addresses)

[amsg(): 100 bytes]
[RA]                   <- current SP
[two amsg arguments]
[something else]
[main() local variables]

(upper addresses)
\end{lstlisting}

amsg()が\main{}に制御を返すとき、ここまでは順調です。
しかし\main{}から\printf{}が呼び出され、これは自身のニーズのためにスタックを使用し、100バイトバッファを破壊します。
良くても、ランダムなゴミが印刷されます。

信じがたいことですが、この問題を修正する方法を知っています：

\begin{lstlisting}[style=customc]
#include <stdio.h>

char* amsg(int n, char* s)
{
        char buf[100];

        sprintf (buf, "error %d: %s\n", n, s) ;

        return buf;
};

char* interim (int n, char* s)
{
        char large_buf[8000];
        // make use of local array.
        // it will be optimized away otherwise, as useless.
        large_buf[0]=0;
        return amsg (n, s);
};

int main()
{
        printf ("%s\n", interim (1234, "something wrong!"));
};
\end{lstlisting}

最適化なしで、\TT{/GS-}オプション\footnote{バッファセキュリティチェックを無効にする}を付けてMSVC 2013でコンパイルすると動作します。
MSVCは警告します：``warning C4172: returning address of local variable or temporary''、しかしコードは実行され、メッセージが印刷されます。
amsg()がinterim()に制御を返す瞬間のスタック状態を見てみましょう：

\begin{lstlisting}
(lower addresses)

[amsg(): 100 bytes]
[RA]                                     <- current SP
[two amsg() arguments]
[interim() stuff, incl. 8000 bytes]
[something else]
[main() local variables]

(upper addresses)
\end{lstlisting}

interim()が\main{}に制御を返す瞬間のスタック状態：

\begin{lstlisting}
(lower addresses)

[amsg(): 100 bytes]
[RA]
[two amsg() arguments]
[interim() stuff, incl. 8000 bytes]
[something else]                         <- current SP
[main() local variables]

(upper addresses)
\end{lstlisting}

\main{}が\printf{}を呼び出すとき、interim()のバッファが割り当てられていた場所でスタックを使用し、
内部にエラーメッセージを持つ100バイトを破壊しません。なぜなら8000バイト（またはそれより少ない）が\printf{}や他の下位関数が行うすべてに十分だからです！

次のような多くの関数がある場合にも動作するかもしれません：
\main{} $\rightarrow$ f1() $\rightarrow$ f2() $\rightarrow$ f3() \dots{} $\rightarrow$ amsg()、
そしてamsg()の結果が\main{}で使用される場合。
\main{}の\ac{SP}と\TT{buf[]}のアドレスの間の距離は十分長くなければなりません。

これが、このようなバグが危険な理由です：時々コードは動作し（バグが気付かれずに隠れる可能性があります）、時々動作しません。
\label{heisenbug}
\myindex{Heisenbug}
このようなバグは冗談でハイゼンバグまたはシュレディンバグ\footnote{\url{https://en.wikipedia.org/wiki/Heisenbug}}と呼ばれます。